\chapter{1977年江苏盐城地区初试卷}

\begin{problem}
计算：

\begin{enumerate}
\item \( (-1)^3-\left[\left(-\frac{1}{2}\right)^2+4^{-\frac{1}{2}}\right]\div\left(\frac{2}{3}\right)^{-2}+\left(\frac{2}{9}\right)^0 \).

\item \( \log_2\frac{1}{2}+\log_3\sqrt{3}+\lg 0.01+\lg 100 \).
\end{enumerate}
\end{problem}

\begin{solution}
\begin{enumerate}
\item 因为 \(4^{-\frac{1}{2}}=\frac{1}{\sqrt{4}}=\frac{1}{2}\)，\(\left(\frac{2}{9}\right)^0=1\)，且 \(\left(\frac{2}{3}\right)^{-2}=\left(\frac{3}{2}\right)^2=\frac{9}{4}\)，所以
\[
(-1)^3-\left[\left(-\frac{1}{2}\right)^2+4^{-\frac{1}{2}}\right]\div\left(\frac{2}{3}\right)^{-2}+\left(\frac{2}{9}\right)^0
=-1-\left(\frac{1}{4}+\frac{1}{2}\right)\div\frac{9}{4}+1
=-1-\frac{3}{4}\times\frac{4}{9}+1=-\frac{1}{3}
\]

\item \(\log_2\frac{1}{2}=-1\)，\(\log_3\sqrt{3}=\frac{1}{2}\)，\(\lg0.01=\lg10^{-2}=-2\)，\(\lg100=2\)，所以原式
\[
=-1+\frac{1}{2}-2+2=-\frac{1}{2}
\]
\end{enumerate}
\end{solution}

\begin{problem}
已知函数 \( y=\sqrt{12-x}+\sqrt{2x+10} \).

\begin{enumerate}
\item 求自变量 \(x\) 的取值范围.

\item 求当 \(x=4\) 时函数 \(y\) 的值.
\end{enumerate}
\end{problem}

\begin{solution}
\begin{enumerate}
\item 要使函数有意义，必须有 \(12-x\geqslant0\) 且 \(2x+10\geqslant0\)，即 \(x\leqslant12\) 且 \(x\geqslant-5\)，所以自变量 \(x\) 的取值范围是 \([-5,12]\).

\item 当 \(x=4\) 时，\(y=\sqrt{8}+\sqrt{18}=2\sqrt{2}+3\sqrt{2}=5\sqrt{2}\).
\end{enumerate}
\end{solution}

\begin{problem}
\(m\) 为何值时，方程 \(x^2-2x+m=0\) 有两个相等的实数根？并求出它的根.
\end{problem}

\begin{solution}
二次方程有两个相等的实数根，当且仅当判别式为零. 因此
\[
\Delta=(-2)^2-4m=4-4m=0
\]
解得 \(m=1\). 此时方程为 \(x^2-2x+1=0\)，即 \((x-1)^2=0\)，所以它的两个相等的实数根都是 \(x=1\).
\end{solution}

\begin{problem}
从 \(\odot O\) 外一点 \(A\) 作 \(\odot O\) 的两条切线 \(AC\)，\(AD\)，\(C\)，\(D\) 为切点，连接 \(CO\) 并延长交 \(AD\) 的延长线于 \(B\). 若 \(AD=4\mathrm{~cm}\)，\(DB=5\mathrm{~cm}\)，求 \(\odot O\) 的面积（\(\pi\) 不必取近似值）.

\FigureLayoutDeclare{img/1977/jiangsu_yancheng_region/q4_fig1.png}{10.0cm}{0bp 26bp 24bp 19bp}
\bitmapfigure[float=inline,max width=6.2cm]{img/1977/jiangsu_yancheng_region/q4_fig1.png}
\FloatBarrier
\end{problem}

\begin{solution}
设圆的半径为 \(r\). 以切点 \(D\) 为原点，直线 \(AD\) 为 \(x\) 轴建立直角坐标系，则由半径垂直于切线可设 \(D=(0,0)\)，\(O=(0,r)\)，\(A=(-4,0)\)，\(B=(5,0)\).
设 \(C=(u,v)\). 因为 \(C\) 在圆上，所以
\[
u^2+(v-r)^2=r^2\quad\text{即}\quad u^2+v^2-2rv=0
\]
又因为 \(AC\) 是圆的切线，故 \(AC\perp OC\)，从而
\[
(-4-u,-v)\cdot(u,v-r)=0
\]
结合 \(u^2+v^2=2rv\)，上式化为 \(4u+rv=0\)，即 \(u=-\frac{rv}{4}\). 代入圆的方程，并注意 \(C\neq D\)，得
\[
v\left(\frac{r^2+16}{16}v-2r\right)=0
\]
因 \(C\neq D\)，有 \(v\neq0\)，故 \(v=\frac{32r}{r^2+16}\)，再由 \(u=-\frac{rv}{4}\) 得 \(u=-\frac{8r^2}{r^2+16}\).
由于 \(B\)，\(O\)，\(C\) 三点共线，直线 \(BO\) 的斜率为 \(-\frac{r}{5}\)，所以
\[
v=-\frac{r}{5}(u-5)=\frac{r}{5}(5-u)
\]
代入 \(u\)，\(v\) 的表达式并约去 \(r>0\)，可得
\[
\frac{32}{r^2+16}=\frac{1}{5}\left(5+\frac{8r^2}{r^2+16}\right)
\]
即 \(160=5(r^2+16)+8r^2=13r^2+80\)，因此 \(r^2=\frac{80}{13}\). 所以圆的面积为
\[
\pi r^2=\frac{80\pi}{13}\mathrm{~cm}^2
\]
\end{solution}

\begin{problem}
求以双曲线 \(3x^2-2y^2=6\) 的顶点为焦点，焦点为顶点的椭圆的方程.
\end{problem}

\begin{solution}
将双曲线方程化为标准形式，得 \(\frac{x^2}{2}-\frac{y^2}{3}=1\). 因此双曲线的顶点为 \((\pm\sqrt{2},0)\)，焦点为 \((\pm\sqrt{5},0)\)，因为其焦半距的平方为 \(2+3=5\).

所求椭圆以双曲线的焦点为顶点，所以椭圆的长半轴为 \(a=\sqrt{5}\)；以双曲线的顶点为焦点，所以椭圆的焦半距为 \(c=\sqrt{2}\). 对椭圆有 \(b^2=a^2-c^2\)，从而 \(b^2=5-2=3\). 因此所求椭圆的方程为
\[
\frac{x^2}{5}+\frac{y^2}{3}=1
\]
\end{solution}

\begin{problem}
若 \(\sin\alpha\)，\(\cos\alpha\)，\(\tan\alpha\) 成等比数列，求证：\(\left(\sin\alpha+\cos\alpha+\tan\alpha\right)\left(\sin\alpha-\cos\alpha+\tan\alpha\right)=\frac{1}{\cos^2\alpha} \).
\end{problem}

\begin{solution}
因为 \(\tan\alpha\) 有定义，所以 \(\cos\alpha\neq0\). 三个数 \(\sin\alpha\)，\(\cos\alpha\)，\(\tan\alpha\) 成等比数列，故中间项的平方等于两端项的积，即
\[
\cos^2\alpha=\sin\alpha\tan\alpha
\]
于是
\[
\begin{gathered}
\left(\sin\alpha+\cos\alpha+\tan\alpha\right)\left(\sin\alpha-\cos\alpha+\tan\alpha\right)\\
=\left(\sin\alpha+\tan\alpha\right)^2-\cos^2\alpha
=\sin^2\alpha+2\sin\alpha\tan\alpha+\tan^2\alpha-\cos^2\alpha\\
=\sin^2\alpha+\cos^2\alpha+\tan^2\alpha
=1+\tan^2\alpha=\frac{1}{\cos^2\alpha}.
\end{gathered}
\]
所以原等式成立.
\end{solution}

\begin{problem}
在边长为 \(1\) 的正方形 \(ABCD\) 内，作内接正方形 \(A_1B_1C_1D_1\).

\begin{enumerate}
\item 在什么情况下，正方形 \(A_1B_1C_1D_1\) 的面积最小？最小值是多少？

\item 在正方形 \(A_1B_1C_1D_1\) 内，再作面积最小的内接正方形 \(A_2B_2C_2D_2\)，\(\ldots\)，如果继续下去，求第 \(n\) 个面积最小的内接正方形的面积.

\item 这样无限制地作下去，可以得到一系列的内接正方形，求它们的面积之和.
\end{enumerate}

\bitmapfigure[max width=6.2cm]{img/1977/jiangsu_yancheng_region/q7_fig1.png}
\end{problem}

\begin{solution}
\begin{enumerate}
\item 设正方形 \(ABCD\) 的顶点坐标为 \(B=(0,0)\)，\(C=(1,0)\)，\(D=(1,1)\)，\(A=(0,1)\). 设 \(A_1=(0,x)\)，\(B_1=(y,0)\)，其中 \(0\leqslant x\)，\(y\leqslant1\). 则
\[
\overrightarrow{A_1B_1}=(y,-x)
\]
由于 \(A_1B_1C_1D_1\) 是正方形，\(\overrightarrow{B_1C_1}\) 是 \(\overrightarrow{A_1B_1}\) 逆时针旋转 \(90^\circ\) 所得，故 \(\overrightarrow{B_1C_1}=(x,y)\). 又因为 \(C_1\) 在直线 \(CD\) 上，其横坐标为 \(1\)，所以 \(x+y=1\).

内接正方形的面积等于边长的平方，因此
\[
S_1=A_1B_1^2=x^2+y^2=\frac{(x+y)^2+(x-y)^2}{2}=\frac{1+(x-y)^2}{2}\geqslant\frac{1}{2}
\]
当且仅当 \(x=y=\frac{1}{2}\) 时取等号，此时四个顶点分别是外正方形四边的中点，且最小面积为 \(\frac{1}{2}\).

\item 设某一个正方形的面积为 \(S\). 由第（1）问，对它作出的面积最小的内接正方形面积为 \(\frac{S}{2}\). 外正方形 \(ABCD\) 的面积为 \(1\)，故第 \(n\) 个面积最小的内接正方形面积为
\(S_n=\left(\frac{1}{2}\right)^n\)，\(n\geqslant1\).

\item 这些正方形的面积构成首项为 \(\frac{1}{2}\)、公比为 \(\frac{1}{2}\) 的无穷等比数列，所以面积之和为
\[
\sum_{n=1}^{\infty}S_n=\frac{\frac{1}{2}}{1-\frac{1}{2}}=1
\]
\end{enumerate}
\end{solution}

\begin{problem}
在 \(\triangle ABC\) 中，最大角是最小角 \(C\) 的二倍，其三边 \(a\)，\(b\)，\(c\) 成等差数列，求证：这个三角形是锐角三角形.
\end{problem}

\begin{solution}
设最小角 \(C=\gamma\). 由题意，最大角为 \(2\gamma\)，第三个角为 \(\pi-3\gamma\). 因为 \(\gamma\) 是最小角，且 \(2\gamma\) 是最大角，所以
\[
\gamma\leqslant\pi-3\gamma\leqslant2\gamma
\]
由左边不等式得 \(\gamma\leqslant\frac{\pi}{4}\).

若 \(\gamma=\frac{\pi}{4}\)，则三角形的三个角为 \(\frac{\pi}{4}\)，\(\frac{\pi}{4}\)，\(\frac{\pi}{2}\). 由正弦定理，与前两个相等角对应的两边相等，而与 \(\frac{\pi}{2}\) 对应的边较长，因此三边中有两个相等且另一个较大. 三个数成等差数列时，若其中两个数相等，则公差必须为零，三个数应全部相等，矛盾. 所以 \(\gamma<\frac{\pi}{4}\).

因此最大角 \(2\gamma<\frac{\pi}{2}\)，三角形的三个内角都小于 \(\frac{\pi}{2}\)，所以这个三角形是锐角三角形.
\end{solution}
